% ============================================
% Supplementary Note 4: Pairwise Selection Correlation
% ============================================
\subsection*{Supplementary Note 4: Pairwise Selection Correlation}
\label{sec:suppl:note4}

To quantify whether species from the same parental community exhibit correlated selection during coalescence, we developed a pairwise selection correlation metric. For each coalescence event, species were assigned to their parental community of origin (the parent with higher abundance if present in both). For each species pair, we evaluated whether their presence/absence patterns in the offspring were concordant (both present or both absent) or discordant (one survives, one extinct).

The concordance rate was converted to a correlation metric: $\rho = 2 \times \text{Concordance rate} - 1$, yielding values from $-1$ (all discordant) to $+1$ (all concordant). We computed $\rho_{\text{same}}$ for within-community pairs and $\rho_{\text{cross}}$ for cross-community pairs, then averaged across coalescence events. To test significance, we generated null distributions by shuffling species origin labels (1,000 permutations) and computed permutation $p$-values for $\Delta = \bar{\rho}_{\text{same}} - \bar{\rho}_{\text{cross}}$.

Positive $\bar{\rho}_{\text{same}}$ indicates species from the same parental community share fates (survive or go extinct together); negative $\bar{\rho}_{\text{cross}}$ indicates cross-community species have opposite fates. In both simulations (Extended Data Fig.~5a--d) and experiments (Extended Data Fig.~6), we observed $\bar{\rho}_{\text{same}} > 0$ and $\bar{\rho}_{\text{cross}} < 0$, with $\Delta$ increasing with interaction strength. At weak interactions ($\mu = 0.3$), both within-community and cross-community pairs show positive correlations with only a small difference ($\Delta = 0.05$; Extended Data Fig.~5a). As interaction strength increases, the difference becomes pronounced: at $\mu = 0.6$, within-community pairs show positive correlation while cross-community pairs become negative ($\Delta = 0.58$; Extended Data Fig.~5b); at $\mu = 0.8$, the separation is even larger ($\Delta = 0.94$; Extended Data Fig.~5c). This confirms that assembly history creates positive pairwise selection correlation among within-community species, providing the mechanistic basis for community-level selection.

\subsection*{Connection to invasion fitness}
\label{sec:suppl:note4_invasion_fitness}

\rev{The pairwise selection correlation can be re-read through the lens of invasion fitness in the generalised Lotka--Volterra (gLV) framework. Invasion fitness is defined as the per-capita growth rate of a species in the rare-invader limit when introduced into a resident community sitting at its stable equilibrium; positive invasion fitness implies successful establishment, negative implies exclusion. This provides a two-level mapping to our metrics: pairwise invasion assays approximate two-species invasion fitness between isolate pairs using a 5\% initial invader frequency, while the pairwise selection correlation $\bar{\rho}_{\text{same}} - \bar{\rho}_{\text{cross}}$ measures correlated multi-species invasion success when a parental community confronts a resident assembled from a different parent. As an auxiliary gLV analysis focused on the within-parent component of this metric, we computed excess same-parent selection correlation, defined as observed same-parent concordance minus an origin-shuffled null that randomises parent-of-origin labels while preserving the coalesced community composition. Across the full simulated $\mu$ sweep ($\mu = 0.05$--$1.20$), this same-parent excess tracks the mean interaction strength with Pearson $r = 0.870$ ($p = 3.2 \times 10^{-8}$; Supplementary Fig.~28). Because the null holds composition fixed and only permutes parent-of-origin labels, this enrichment is not expected from the label structure alone; it is consistent with interaction-mediated correlated invasion outcomes.}
